\label{chap:introduction}

A webalkalmazások napjainkban a digitális szolgáltatások egyik legelterjedtebb formáját jelentik. Böngészőből elérhetők, több eszközön használhatók, és gyakran egyszerre kell megfelelniük felhasználói élményre, teljesítményre, biztonságra, skálázhatóságra és karbantarthatóságra vonatkozó elvárásoknak. Egy modern webalkalmazás fejlesztése ezért nem csupán programozási feladat, hanem tervezési, termékfejlesztési és technológiai döntések sorozata is.

A dolgozat célja egy modern webalkalmazás tervezési és fejlesztési folyamatának bemutatása egy konkrét projekt, a WatchDB példáján keresztül. A hangsúly nem kizárólag az elkészült alkalmazáson van, hanem azon a folyamaton is, amely a projektötlettől a követelmények meghatározásán, az MVP (Minimum Viable Product) kijelölésén, a technológiai döntéseken és az implementáción keresztül vezet a működő rendszerig. Ennek megfelelően az alkalmazás a dolgozatban esettanulmányként jelenik meg: segítségével bemutatható, hogyan válhat egy termékkoncepció strukturált fejlesztési feladattá.

A WatchDB egy filmkövető webalkalmazás koncepciójára épül. A projekt kiinduló problémája, hogy a felhasználók különböző platformokon néznek filmeket, miközben gyakran nincs egységes helyük arra, hogy nyilvántartsák a megnézett filmeket, vezessék a később megtekintendő filmek listáját, illetve értékeljék saját filmélményeiket. A termék hosszabb távú víziója közösségi irányba is bővíthető: a felhasználók követhetik egymás aktivitását, felfedezhetnek ajánlásokat, és megoszthatják filmnézési szokásaikat. A dolgozatban bemutatott fejlesztési fókusz ugyanakkor az első, megvalósítható változat, vagyis az MVP köré szerveződik.

A projekt előkészítésének fontos eleme a Product Requirements Document (PRD), amely a termék célját, célközönségét, fő funkcióit, sikerességi szempontjait, korlátait és fejlesztési ütemezését foglalja össze. A PRD szerepe ebben a dolgozatban az, hogy láthatóvá tegye: a fejlesztés nem közvetlenül kódolással indul, hanem a probléma, a felhasználói igények és a megvalósítható terjedelem tisztázásával. Ez különösen fontos az MVP meghatározásánál, ahol el kell dönteni, hogy a teljes termékvízióból mely funkciók szükségesek az első használható verzióhoz, és melyek kerülhetnek későbbi fejlesztési szakaszba.

A dolgozat felépítése ezt a gondolatmenetet követi. Az irodalmi áttekintés bemutatja a webalkalmazások fejlődését, a modern webfejlesztés alapfogalmait, a fejlesztési módszertanokat, a webfejlesztői csapat szerepköreit, az SDLC folyamatát és a választott technológiai stack szakmai hátterét. Ezt követően a projektcélok és termékkoncepció fejezete a WatchDB PRD-jéből kiindulva ismerteti a termék célját és tervezési keretét. A további fejezetek a követelmények, architekturális döntések, gyakorlati megvalósítás és értékelés bemutatásán keresztül tárgyalják, hogyan jelennek meg ezek a döntések egy konkrét webalkalmazás fejlesztése során.
